\documentclass[11pt]{article}

\usepackage[margin=1in]{geometry}
\usepackage{amsmath}
\usepackage{amssymb}
\usepackage{booktabs}
\usepackage{graphicx}
\usepackage{xcolor}
\usepackage{listings}
\usepackage{enumitem}
\usepackage{caption}
\usepackage[hidelinks]{hyperref}
\usepackage{siunitx}
% Preserve the manual's comma digit grouping (e.g. 101,000) for pressures.
\sisetup{group-separator={,}, group-minimum-digits=4}

% ---------------------------------------------------------------------------
% Listings configuration for Python and shell snippets
% ---------------------------------------------------------------------------
\definecolor{codebg}{rgb}{0.97,0.97,0.97}
\definecolor{codekw}{rgb}{0.0,0.2,0.7}
\definecolor{codecomment}{rgb}{0.0,0.5,0.0}
\definecolor{codestring}{rgb}{0.6,0.1,0.1}

\lstdefinestyle{pystyle}{
  language=Python,
  backgroundcolor=\color{codebg},
  basicstyle=\ttfamily\small,
  keywordstyle=\color{codekw}\bfseries,
  commentstyle=\color{codecomment}\itshape,
  stringstyle=\color{codestring},
  numbers=left,
  numberstyle=\tiny\color{gray},
  frame=single,
  rulecolor=\color{gray!40},
  breaklines=true,
  showstringspaces=false,
  tabsize=4,
}

\lstdefinestyle{shellstyle}{
  backgroundcolor=\color{codebg},
  basicstyle=\ttfamily\small,
  frame=single,
  rulecolor=\color{gray!40},
  breaklines=true,
  showstringspaces=false,
}

% Inline monospace that is allowed to break across lines (so long option
% names, file names and identifiers do not overrun the right margin). Breaks
% are permitted after underscores in addition to TeX's default break after
% explicit hyphens.
\newcommand{\codeus}{\textunderscore\allowbreak}
\DeclareRobustCommand{\code}[1]{\texttt{\ifmmode\else\renewcommand{\_}{\codeus}\fi#1}}

% Allow TeX a little extra stretch on the final pass so prose lines that end
% with a long code token or identifier do not overrun the right margin.
\setlength{\emergencystretch}{3em}

\title{\textbf{WindowBending}\\[4pt]
\large Nonlinear Deformation Analysis of Pressurized Circular Windows\\[6pt]
\normalsize User Manual}
\author{}
\date{\today}

\begin{document}

\maketitle
\tableofcontents
\newpage

% ===========================================================================
\section{Introduction}
% ===========================================================================

\textbf{WindowBending} is a Python tool that computes the structural response
of a circular window (for example, an optical fused-silica window) subjected to
a uniform pressure load on one face. It solves the axisymmetric, large-deflection
F\"oppl--von K\'arm\'an plate equations and reports:

\begin{itemize}[noitemsep]
  \item the transverse deflection profile $w(r)$ and the centre deflection;
  \item the radial and circumferential (hoop) stress distributions on both faces;
  \item the peak tensile stress; and
  \item the safety factor against a user-supplied allowable stress.
\end{itemize}

The tool supports two edge support models (\emph{clamped} and
\emph{simply supported}) and an optional \emph{design sweep} that varies one
parameter (thickness or pressure) over a range and tabulates the resulting
deflection, stress, and safety factor.

\paragraph{Motivating problem.} The default configuration models a
\qty{450}{\milli\meter} diameter, \qty{25}{\milli\meter} thick fused-silica window,
fixed around its outer edge, with a uniform pressure of
\qty{101000}{\pascal} (1~atm) applied to the top face.

% ===========================================================================
\section{Theory}
% ===========================================================================

\subsection{Geometry and assumptions}

The window is modelled as a thin, flat, circular plate of radius $a$ and uniform
thickness $t$, made of a linear-elastic, isotropic material with Young's modulus
$E$ and Poisson's ratio $\nu$. Because both the geometry and the loading are
axisymmetric, all field quantities depend only on the radial coordinate
$r \in [0, a]$. The transverse deflection is denoted $w(r)$ (positive in the
direction of the applied pressure) and the in-plane radial displacement is
$u(r)$. The \emph{centre deflection} $w_{0} \equiv w(0)$ is the value of this
field at the centre $r = 0$, which is the maximum deflection; the dimensionless
ratio $w_{0}/t$ measures how large the deflection is relative to the plate
thickness and is used throughout as the indicator of large-deflection behaviour.

\subsection{Linear (small-deflection) plate theory}
\label{sec:linear}

For small deflections ($w \ll t$), Kirchhoff--Love plate theory gives a closed-form
centre deflection for a uniformly loaded circular plate. The flexural rigidity is
%
\begin{equation}
  D = \frac{E\,t^{3}}{12\,(1-\nu^{2})}.
\end{equation}
%
For a \textbf{clamped} edge,
%
\begin{equation}
  w_{0}^{\text{lin}} = \frac{p\,a^{4}}{64\,D},
  \label{eq:clamped-linear}
\end{equation}
%
and for a \textbf{simply supported} edge,
%
\begin{equation}
  w_{0}^{\text{lin}} = \frac{p\,a^{4}\,(5+\nu)}{64\,D\,(1+\nu)}.
  \label{eq:ss-linear}
\end{equation}
%
These expressions are used both as a sanity baseline and to seed the nonlinear
solver. They are exposed by \code{PlateConfig.linear\_center\_deflection()}.

\subsection{Large-deflection F\"oppl--von K\'arm\'an equations}

When the deflection becomes a non-negligible fraction of the thickness, the plate
develops in-plane (membrane) stretching that stiffens the response. This coupling
is captured by the axisymmetric F\"oppl--von K\'arm\'an equations. The membrane
stiffness is
%
\begin{equation}
  C = \frac{E\,t}{1-\nu^{2}},
\end{equation}
%
and the radial and circumferential membrane stress resultants are
%
\begin{align}
  N_{r} &= C\left( u' + \nu\,\frac{u}{r} + \tfrac{1}{2}\,(w')^{2} \right), \\
  N_{\theta} &= C\left( \frac{u}{r} + \nu\,u' + \tfrac{1}{2}\,\nu\,(w')^{2} \right),
\end{align}
%
where a prime denotes $\mathrm{d}/\mathrm{d}r$. The coupled governing equations are
the transverse equilibrium equation,
%
\begin{equation}
  D\,\nabla^{4} w = p + N_{r}\,w'' + N_{\theta}\,\frac{w'}{r},
  \label{eq:transverse}
\end{equation}
%
and the in-plane equilibrium equation,
%
\begin{equation}
  \frac{\mathrm{d}}{\mathrm{d}r}\!\left( N_{r} \right) + \frac{N_{r}-N_{\theta}}{r} = 0,
  \label{eq:inplane}
\end{equation}
%
where the axisymmetric biharmonic operator acting on $w$ is
%
\begin{equation}
  \nabla^{4} w
  = w'''' + \frac{2}{r}\,w''' - \frac{1}{r^{2}}\,w'' + \frac{1}{r^{3}}\,w'.
\end{equation}

\subsection{Range of validity}
\label{sec:validity}

The Kirchhoff and F\"oppl--von K\'arm\'an models rest on two independent
assumptions, each with its own dimensionless ratio:
%
\begin{itemize}
  \item \textbf{Thin-plate (shear-free) assumption:} the thickness must be small
    relative to the in-plane span, quantified by the thickness-to-diameter ratio
    $t/D$ (equivalently the radius-to-thickness ratio $R/h = 1/(2\,t/D)$, with
    $R=a$ the radius and $h=t$ the thickness). Kirchhoff--Love kinematics neglect
    transverse shear and are accurate for genuinely thin plates,
    $t/D < 0.05$ ($R/h > 10$). Beyond this the plate carries an appreciable shear
    deflection that Kirchhoff theory omits: a \emph{transition zone}
    $0.05 \le t/D < 0.1$ ($5 < R/h \le 10$) where the Mindlin--Reissner model is
    recommended for accuracy, and a definite thick-plate region $t/D \ge 0.1$
    ($R/h \le 5$) where it is required.
  \item \textbf{Moderate-rotation assumption:} the F\"oppl--von K\'arm\'an strain
    measure retains only the leading nonlinear term $\tfrac{1}{2}(w')^{2}$, which
    is accurate while the central deflection is at most a few thicknesses
    ($w_0/t \lesssim 1$--$5$) and edge rotations stay below about \qtyrange{10}{15}{\degree}.
\end{itemize}
%
For the windows of interest the deflection stays far below the thickness
($w_0/t \le 0.06$), so the moderate-rotation limit is comfortably satisfied and
the response is effectively quasi-linear. The binding limitation is instead the
thin-plate assumption: much of the design space---including the default
\qty{450}{\milli\meter} diameter, \qty{25}{\milli\meter} thick window
($t/D = 0.056$, $R/h = 9$)---falls in
the transition zone or beyond ($t/D$ reaches $0.14$ at the thick, small-diameter
corners). There the shear-free Kirchhoff model under-predicts the deflection,
motivating the shear-deformable model of the next subsection.

\subsection{Mindlin--Reissner (shear-deformable) theory}
\label{sec:mindlin}

When $t/D$ exceeds the thin-plate limit the tool can solve the axisymmetric
\emph{Mindlin--Reissner} plate equations, which relax the Kirchhoff constraint by
letting the cross-section rotation $\beta(r)$ be independent of the slope $w'$.
The difference between them is the transverse-shear strain, resisted by the shear
stiffness
%
\begin{equation}
  k_{s} = \kappa\,G\,t, \qquad
  G = \frac{E}{2(1+\nu)}, \qquad
  \kappa = \tfrac{5}{6},
\end{equation}
%
with $\kappa$ the shear-correction factor for a homogeneous rectangular section.
The transverse shear resultant is $Q = k_{s}\,(w' + \beta)$, and the bending
moments are expressed through $\beta$ rather than $w''$:
%
\begin{align}
  M_{r} &= D\left( \beta' + \nu\,\frac{\beta}{r} \right), &
  M_{\theta} &= D\left( \nu\,\beta' + \frac{\beta}{r} \right).
\end{align}
%
Combined with the same membrane (von K\'arm\'an) coupling used above, this yields
a shear-deformable BVP that reduces to the Kirchhoff solution as $t/D \to 0$
(where $k_{s} \to \infty$ forces $\beta \to -w'$). For a clamped plate the linear
limit gives the closed-form central deflection
%
\begin{equation}
  w_{\mathrm{Mindlin}}(0) = \underbrace{\frac{p\,a^{4}}{64\,D}}_{\text{bending}}
  + \underbrace{\frac{p\,a^{2}}{4\,\kappa\,G\,t}}_{\text{shear}},
  \label{eq:mindlin-clamped}
\end{equation}
%
i.e.\ the Kirchhoff deflection plus an additive shear term. The peak bending
stress is essentially unchanged (transverse equilibrium fixes $Q=-pr/2$ in both
theories); only the deflection grows, by the shear compliance in
\eqref{eq:mindlin-clamped}. The theory is selected with the
\code{-{}-plate-theory} flag (Section~\ref{sec:cli}); \code{auto} switches to
Mindlin whenever $R/h \le 10$ (equivalently $t/D \ge 0.05$), so the transition
zone and thicker plates use the shear-deformable model, and Kirchhoff is retained
only for genuinely thin plates ($R/h > 10$).



Three conditions follow from axisymmetry at the centre $r=0$:
%
\begin{equation}
  w'(0) = 0, \qquad w'''(0) = 0, \qquad u(0) = 0.
\end{equation}
%
At the supported edge $r=a$, the deflection vanishes and the edge is radially
constrained:
%
\begin{equation}
  w(a) = 0, \qquad u(a) = 0.
\end{equation}
%
The remaining condition distinguishes the two support types:
%
\begin{itemize}
  \item \textbf{Clamped:} zero edge rotation,
  \begin{equation}
    w'(a) = 0.
  \end{equation}
  \item \textbf{Simply supported:} zero radial bending moment,
  \begin{equation}
    M_{r}(a) = -D\left( w''(a) + \nu\,\frac{w'(a)}{a} \right) = 0.
  \end{equation}
\end{itemize}

\subsection{Stress recovery}

Once $w(r)$ and $u(r)$ are known, the bending moments per unit length are
%
\begin{align}
  M_{r} &= -D\left( w'' + \nu\,\frac{w'}{r} \right), &
  M_{\theta} &= -D\left( \frac{w'}{r} + \nu\,w'' \right).
\end{align}
%
The total surface stress superposes a membrane part (uniform through thickness)
and a bending part (linear through thickness, extreme at the faces):
%
\begin{align}
  \sigma_{r}^{\text{membrane}} = \frac{N_{r}}{t}, &\qquad
  \sigma_{\theta}^{\text{membrane}} = \frac{N_{\theta}}{t}, \\[4pt]
  \sigma_{r}^{\text{bending}} = \frac{6 M_{r}}{t^{2}}, &\qquad
  \sigma_{\theta}^{\text{bending}} = \frac{6 M_{\theta}}{t^{2}}.
\end{align}
%
The stresses on the two faces are therefore
%
\begin{align}
  \sigma^{\text{top}} &= -\sigma^{\text{bending}} + \sigma^{\text{membrane}}, \\
  \sigma^{\text{bottom}} &= +\sigma^{\text{bending}} + \sigma^{\text{membrane}}.
\end{align}

\subsection{Peak stress and safety factor}

The peak tensile stress is the maximum over the radial coordinate of all four
stress components (radial/hoop on top/bottom faces):
%
\begin{equation}
  \sigma_{\max} = \max_{r}\, \bigl\{\,
    \sigma_{r}^{\text{top}},\ \sigma_{\theta}^{\text{top}},\
    \sigma_{r}^{\text{bottom}},\ \sigma_{\theta}^{\text{bottom}}
  \,\bigr\}.
\end{equation}
%
The safety factor compares this against the user-supplied allowable stress
$\sigma_{\text{allow}}$:
%
\begin{equation}
  \text{SF} = \frac{\sigma_{\text{allow}}}{\sigma_{\max}}.
\end{equation}
%
Because fused silica is brittle, the relevant allowable stress is its
``allowable design stress'' as listed in the reference glass-property table
(Ref.~\ref{ref:bnl}).
The default allowable stress is \qty{4.69}{\mega\pascal} (\num{680} psi).

% ===========================================================================
\section{Numerical Implementation}
% ===========================================================================

\subsection{Reformulation for numerical conditioning}

A direct first-order reduction of the fourth-order biharmonic operator carries
$1/r^{2}$ and $1/r^{3}$ singular terms near the centre. In practice these make the
collocation solver \code{scipy.integrate.solve\_bvp} ill-conditioned: it fails to
converge and returns a deflection biased low by roughly $7\%$, even in the linear
limit where the answer must match Eqs.~\eqref{eq:clamped-linear}
and~\eqref{eq:ss-linear}.

To avoid this, the transverse equation is rewritten using the Laplacian
intermediate variable
%
\begin{equation}
  M \equiv \nabla^{2} w = w'' + \frac{w'}{r}.
\end{equation}
%
This leaves only mild $1/r$ singularities. The simply-supported moment condition
becomes, using $w'' = M - w'/r$,
%
\begin{equation}
  M(a) = (1-\nu)\,\frac{w'(a)}{a}.
\end{equation}
%
With this reformulation the solver reproduces both closed-form linear solutions to
ten significant digits.

\subsection{First-order state-space system}

The problem is cast as a first-order boundary-value problem in the state vector
%
\begin{equation}
  \mathbf{y} = \bigl[\, w,\ w',\ M,\ M',\ u,\ u' \,\bigr]^{\mathsf{T}},
\end{equation}
%
integrated over $r \in [r_{0}, a]$, where $r_{0} = a \times 10^{-6}$ is a small
offset that regularizes the residual $1/r$ terms at the centre (constant
\code{CENTER\_REGULARIZATION}). The six boundary residuals encode the three
symmetry conditions at $r_0$ and the three edge conditions at $a$.

\subsection{Solver settings}

The nonlinear BVP is solved with \code{solve\_bvp}, seeded with the clamped
small-deflection shape from Eq.~\eqref{eq:clamped-linear}. Key defaults are listed
in Table~\ref{tab:solver}.

\begin{table}[h]
  \centering
  \caption{Default solver configuration constants.}
  \label{tab:solver}
  \begin{tabular}{lll}
    \toprule
    Constant & Value & Meaning \\
    \midrule
    \code{DEFAULT\_MESH\_NODES}     & 250        & Initial mesh nodes \\
    \code{DEFAULT\_OUTPUT\_POINTS}  & 400        & Output sample points \\
    \code{SOLVER\_TOLERANCE}        & $10^{-7}$  & Collocation residual tolerance \\
    \code{SOLVER\_MAX\_NODES}       & 50{,}000   & Maximum adaptive mesh nodes \\
    \code{CENTER\_REGULARIZATION}   & $10^{-6}$  & Centre offset $r_{0}/a$ \\
    \code{KIRCHHOFF\_CONTINUATION\_STEPS} & 10   & Load-continuation steps on failure \\
    \code{LINEAR\_FALLBACK\_WT\_RATIO} & 0.2      & $w_0/t$ below which linear fallback is used \\
    \code{VON\_KARMAN\_WT\_WARN\_RATIO} & 5       & $w_0/t$ above which validity is warned \\
    \bottomrule
  \end{tabular}
\end{table}

\subsection{Thin-plate convergence and graceful fallback}
\label{sec:robustness}

As the plate is made thinner, the von K\'arm\'an central deflection grows toward
the order of the thickness ($w_0/t \gtrsim 1$) and the response becomes
membrane-dominated. A sharp bending \emph{boundary layer} then forms at the
clamped edge, whose width shrinks with load. A uniform-mesh collocation solver
refines that layer without bound and exhausts its node budget; empirically the
single-shot solve succeeds for the fused-silica window down to about $t=\qty{6}{\milli\meter}$
($w_0/t\approx0.5$) and fails below that. This is a strong-nonlinearity limit,
not a small-deflection conditioning problem.

The solver therefore uses a three-tier strategy:
%
\begin{enumerate}
  \item \textbf{Single-shot solve.} The nonlinear BVP is attempted once from the
    linear initial guess (the common, fast path).
  \item \textbf{Load continuation.} On failure the pressure is ramped from a
    small fraction to its full value in \code{KIRCHHOFF\_CONTINUATION\_STEPS}
    steps, each converged solution seeding the next. This homotopy reaches
    deeper into the nonlinear regime.
  \item \textbf{Graceful fallback.} If continuation also fails, the response is
    inspected: when $w_0/t \le \code{LINEAR\_FALLBACK\_WT\_RATIO}$ ($0.2$) the
    membrane coupling is negligible, so the accurate closed-form linear
    Kirchhoff solution (Section~\ref{sec:linear}) is returned with a warning.
    Otherwise a descriptive \code{RuntimeError} is raised explaining that the
    case is in the strongly nonlinear, membrane-dominated regime.
\end{enumerate}
%
Independently, whenever a converged solution has $w_0/t >
\code{VON\_KARMAN\_WT\_WARN\_RATIO}$ ($5$) the tool warns that the
F\"oppl--von K\'arm\'an moderate-rotation assumption itself is violated
(Section~\ref{sec:validity}) and a large-rotation shell model is required.
Because such a result is quantitatively unreliable, this breach is surfaced
prominently in \emph{every} output channel so it cannot be missed:
%
\begin{itemize}
  \item a Python \code{UserWarning} is emitted;
  \item the console prints a boxed \emph{MODEL VALIDITY WARNING} banner (before
    and after a single case; sweeps additionally flag each offending row/cell
    and list the affected cases);
  \item the saved figures carry a bold red banner across the bottom
    (\emph{MODEL INVALID: deflection exceeds F\"oppl--von K\'arm\'an validity});
  \item the Excel workbooks gain a red-highlighted validity row, and sweep
    workbooks additionally shade the offending value cells red.
\end{itemize}

% ===========================================================================
\section{Software Architecture}
% ===========================================================================

The implementation lives in a single module, \code{window\_deformation.py}, and is
organised around small, immutable data classes and pure functions so that the
numerical core is independent of the command-line interface and plotting.

\subsection{Data classes}

\begin{description}[style=nextline]
  \item[\code{Material}] Young's modulus, Poisson's ratio, allowable stress, and a
    name. Validates that the modulus and allowable stress are positive and that
    $-1 < \nu < 0.5$.
  \item[\code{PlateConfig}] Geometry (diameter, thickness), pressure, material, and
    boundary condition. Provides derived properties \code{radius\_m},
    \code{flexural\_rigidity}, \code{membrane\_stiffness}, and
    \code{linear\_center\_deflection()}. The convenience constructor
    \code{from\_engineering\_units()} accepts millimetre geometry.
  \item[\code{PlateSolution}] The computed fields: $r$, $w$, the linear and
    nonlinear centre deflections, the four stress arrays, the peak stress, the
    safety factor, and the boundary condition. Exposes \code{center\_deflection\_mm}
    and \code{max\_tensile\_mpa} convenience properties.
  \item[\code{SweepPoint}] One evaluated case of a design sweep: the swept value,
    centre deflection (\unit{\milli\meter}), peak stress (\unit{\mega\pascal}), and safety factor.
  \item[\code{GridPoint}] One evaluated case of a combined grid sweep: thickness
    (\unit{\milli\meter}), diameter (\unit{\milli\meter}), thickness/\allowbreak diameter ratio,
    resolved plate theory, linear and nonlinear centre deflection (\unit{\milli\meter}),
    peak stress (\unit{\mega\pascal}), and safety factor.
\end{description}

\subsection{Key functions}

\begin{description}[style=nextline]
  \item[\code{solve\_plate(config, n\_mesh, n\_points)}] Solves the BVP for a single
    \code{PlateConfig} and returns a \code{PlateSolution}.
  \item[\code{run\_design\_sweep(base\_config, sweep\_variable, sweep\_values)}]
    Evaluates the plate response across a range of thickness or pressure and
    returns a list of \code{SweepPoint}.
  \item[\code{run\_sweep\_grid(base\_config, thickness\_values\_mm, diameter\_values\_mm)}]
    Evaluates the response over a full thickness $\times$ diameter grid (pressure,
    material and edge fixed) and returns a list of \code{GridPoint}.
  \item[\code{build\_parser()}, \code{config\_from\_args(args)}] Construct the CLI
    parser and translate parsed arguments into a \code{PlateConfig}.
  \item[\code{main(argv)}] Entry point: dispatches to single-case or sweep mode,
    prints results, and optionally plots.
\end{description}

Plotting helpers import \code{matplotlib} lazily and force the non-interactive Agg
backend, so the numerical core has no hard GUI dependency, remains fast to import
for testing, and figures are written to disk rather than displayed. JPEG output
additionally uses Pillow, which ships as a \code{matplotlib} dependency.

% ===========================================================================
\section{Installation}
% ===========================================================================

The tool requires Python 3.10+ and the packages \code{numpy}, \code{scipy},
\code{matplotlib}, and \code{openpyxl} (for Excel export).

\begin{lstlisting}[style=shellstyle]
python -m pip install numpy scipy matplotlib openpyxl
\end{lstlisting}

% ===========================================================================
\section{Command-Line Usage}
\label{sec:cli}
% ===========================================================================

The script is run directly with Python. All parameters have defaults matching the
fused-silica window described in the introduction.

\begin{lstlisting}[style=shellstyle]
python window_deformation.py [options]
\end{lstlisting}

\subsection{Options}

\begin{table}[h]
  \centering
  \caption{Command-line options.}
  \label{tab:cli}
  \begin{tabular}{@{}l l p{8.2cm}@{}}
    \toprule
    Option & Default & Description \\
    \midrule
    \code{-{}-diameter-mm}         & 450     & Window diameter (\unit{\milli\meter}) \\
    \code{-{}-thickness-mm}        & 25      & Window thickness (\unit{\milli\meter}) \\
    \code{-{}-pressure-pa}         & 101000  & Uniform pressure (\unit{\pascal}) \\
    \code{-{}-material}            & fused\_silica & Named material preset (see Table~\ref{tab:materials}) \\
    \code{-{}-list-materials}      & off     & Print the preset table and exit \\
    \code{-{}-youngs-modulus-pa}   & preset  & Override the preset Young's modulus (\unit{\pascal}) \\
    \code{-{}-poisson-ratio}       & preset  & Override the preset Poisson's ratio \\
    \code{-{}-allowable-stress-mpa}& preset  & Override the preset allowable stress (\unit{\mega\pascal}) \\
    \code{-{}-boundary-condition}  & clamped & \code{clamped} or \code{simply\_supported} \\
    \code{-{}-plate-theory}        & auto    & \code{kirchhoff}, \code{mindlin}, or \code{auto} (Mindlin when $R/h\le10$, i.e.\ $t/D\ge0.05$) \\
    \code{-{}-sweep-variable}      & none    & \code{none}, \code{thickness\_mm}, \code{pressure\_pa}, \code{diameter\_mm}, or \code{all} \\
    \code{-{}-sweep-start}         & ---     & Sweep start value (single-variable sweep only) \\
    \code{-{}-sweep-stop}          & ---     & Sweep stop value (single-variable sweep only) \\
    \code{-{}-sweep-count}         & 7       & Number of sweep samples (single-variable sweep only) \\
    \code{-{}-thickness-sweep-start} & 15    & Thickness-leg start (\unit{\milli\meter}) for the \code{all} grid \\
    \code{-{}-thickness-sweep-stop}  & 35    & Thickness-leg stop (\unit{\milli\meter}) for the \code{all} grid \\
    \code{-{}-thickness-sweep-count} & 5     & Number of thickness points for the \code{all} grid \\
    \code{-{}-diameter-sweep-start}  & 250   & Diameter-leg start (\unit{\milli\meter}) for the \code{all} grid \\
    \code{-{}-diameter-sweep-stop}   & 450   & Diameter-leg stop (\unit{\milli\meter}) for the \code{all} grid \\
    \code{-{}-diameter-sweep-count}  & 5     & Number of diameter points for the \code{all} grid \\
    \code{-{}-output-dir}          & figures & Directory for saved figures and spreadsheets \\
    \code{-{}-figure-name}         & auto    & Base name for figure files (boundary suffix still appended) \\
    \code{-{}-excel-name}          & auto    & Base name for the Excel workbook (boundary suffix still appended) \\
    \code{-{}-figure-format}       & jpg     & One or more of \code{jpg}, \code{png}, \code{svg}, \code{pdf} (e.g.\ \code{jpg svg}) \\
    \code{-{}-no-save}             & off     & Skip writing all output files (figures and spreadsheets) \\
    \code{-{}-no-figures}          & off     & Skip only the figures (still write spreadsheets) \\
    \code{-{}-no-excel}            & off     & Skip only the Excel spreadsheets (still write figures) \\
    \bottomrule
  \end{tabular}
\end{table}

\subsection{Material presets}
\label{sec:materials}

The \code{-{}-material} flag selects a named material preset that sets Young's
modulus $E$, Poisson's ratio $\nu$ and the allowable stress in one step. The
built-in presets are listed in Table~\ref{tab:materials}. Run
\code{-{}-list-materials} to print the same table from the command line and exit.

\begin{table}[h]
  \centering
  \caption{Built-in material presets. Values are nominal room-temperature
  properties taken from Ref.~\ref{ref:bnl}. For the glasses the allowable stress
  is the ``allowable design
  stress'' listed directly in the reference glass-property table; for methyl
  methacrylate (\code{acrylic}) it follows table Note~\#8 (design stress limited
  to the lower of modulus of rupture / 10 or $(9200/10)\times(E/\num{360000})$),
  giving \num{920} psi at \qty{75}{\degree F}. Because option~(b) of that rule
  scales with the elastic modulus, overriding \code{-{}-youngs-modulus-pa} for
  \code{acrylic} recomputes the worst case of the two options: a modulus below
  the \num{360000} psi reference lowers the allowable below \num{920} psi, while
  a modulus at or above it stays capped at \num{920} psi. Verify against
  certified data before relying on them for design.}
  \label{tab:materials}
  \begin{tabular}{llS[table-format=3.1]S[table-format=1.3]S[table-format=3.2]}
    \toprule
    Key & Name & {$E$ (\unit{\giga\pascal})} & {$\nu$} & {Allowable (\unit{\mega\pascal})} \\
    \midrule
    \code{fused\_silica}    & Fused Silica            & 72.0  & 0.170 & 4.69   \\
    \code{n\_bk7}           & N-BK7                   & 82.0  & 0.206 & 3.45   \\
    \code{borofloat}        & Borofloat 33            & 64.0  & 0.200 & 25.00  \\
    \code{aluminosilicate}  & Aluminosilicate (1723)  & 85.5  & 0.260 & 4.55   \\
    \code{vycor}            & 96\% Silica (Vycor)     & 68.9  & 0.190 & 4.24   \\
    \code{pyrex}            & Borosilicate (Pyrex)    & 62.7  & 0.200 & 4.21   \\
    \code{plate\_glass}     & Plate Glass (Herculite) & 68.9  & 0.210 & 4.48   \\
    \code{lead\_glass}      & High-Lead X-Ray Glass   & 55.2  & 0.230 & 3.45   \\
    \code{acrylic}          & Methyl Methacrylate     & 2.5   & 0.390 & 6.34   \\
    \code{sapphire}         & Sapphire                & 345.0 & 0.290 & 350.00 \\
    \code{zerodur}          & Zerodur                 & 90.3  & 0.243 & 57.00  \\
    \code{caf2}             & Calcium Fluoride        & 75.8  & 0.260 & 36.00  \\
    \code{beryllium}        & Beryllium               & 287.0 & 0.032 & 240.00 \\
    \bottomrule
  \end{tabular}
\end{table}

\noindent The individual \code{-{}-youngs-modulus-pa}, \code{-{}-poisson-ratio}
and \code{-{}-allowable-stress-mpa} flags override the corresponding preset value,
so a preset can be used as a starting point and then tuned. For example,
\code{-{}-material sapphire -{}-allowable-stress-mpa 200} keeps sapphire's stiffness
but applies a more conservative allowable stress. The selected material name is
echoed in the console header, the figure subtitles and the spreadsheet parameter
rows.

\paragraph{How to use the preset library.} The presets are used in four steps:
\begin{enumerate}[leftmargin=2em]
  \item \textbf{Discover.} List the available presets and their properties without
    running a solve:
    \begin{lstlisting}[style=shellstyle]
python window_deformation.py --list-materials
    \end{lstlisting}
  \item \textbf{Select.} Pass a preset key (the left column of
    Table~\ref{tab:materials}) to \code{-{}-material}. All other inputs keep their
    defaults unless you set them:
    \begin{lstlisting}[style=shellstyle]
python window_deformation.py --material beryllium --thickness-mm 15
    \end{lstlisting}
  \item \textbf{Tune (optional).} Override any single property while keeping the
    rest of the preset. The override always wins over the preset value:
    \begin{lstlisting}[style=shellstyle]
python window_deformation.py --material n_bk7 --allowable-stress-mpa 25
    \end{lstlisting}
  \item \textbf{Combine.} A preset applies to every case of a design sweep, because
    the sweep holds the material fixed while varying geometry or load:
    \begin{lstlisting}[style=shellstyle]
python window_deformation.py --material sapphire \
    --sweep-variable thickness_mm --sweep-start 10 --sweep-stop 30 --sweep-count 5
    \end{lstlisting}
\end{enumerate}

\noindent If \code{-{}-material} is omitted the tool uses \code{fused\_silica}, so
existing commands behave exactly as before. Supplying an unknown key produces an
\code{argparse} error that lists the valid choices.

\paragraph{Presets from the Python API.} The same presets are available
programmatically through the \code{wd.MATERIAL\_PRESETS} dictionary, keyed by the
identical names accepted by \code{-{}-material}:
\begin{lstlisting}[style=pystyle]
import window_deformation as wd

sapphire = wd.MATERIAL_PRESETS["sapphire"]         # a ready-made Material
config = wd.PlateConfig.from_engineering_units(
    diameter_mm=450.0, thickness_mm=25.0, pressure_pa=101_000.0,
    material=sapphire,
)
solution = wd.solve_plate(config)
print(solution.center_deflection_mm, solution.max_tensile_mpa)
\end{lstlisting}

\noindent Because \code{Material} is an immutable (frozen) dataclass, tune a preset
by constructing a copy with \code{dataclasses.replace} rather than mutating it:
\begin{lstlisting}[style=pystyle]
from dataclasses import replace

conservative = replace(wd.MATERIAL_PRESETS["beryllium"],
                       allowable_stress_pa=200 * wd.PA_PER_MPA)
\end{lstlisting}

\paragraph{Figure output.} The tool never opens an interactive window. Instead it
renders results with the non-interactive Agg backend and writes image files
into the output directory (\code{figures/} by default), reporting each saved path
on the console. A single case produces
\code{single\_case\_deflection\_fused\_silica\_clamped.jpg}
and \code{single\_case\_stress\_fused\_silica\_clamped.jpg}; a single-parameter
sweep produces \code{sweep\_\textit{variable}\_fused\_silica\_clamped.jpg}; the
combined sweep produces \code{sweep\_combined\_fused\_silica\_clamped.jpg} (the
material name and trailing boundary-condition name are appended
automatically---see \emph{Output file naming} below). Every figure is also
annotated with the analysed material and the fixed simulation parameters via a
two-line subtitle (see \emph{Fixed-parameter annotations} below). The figure format
defaults to JPEG but is selectable with \code{-{}-figure-format}, which accepts one
or more of \code{jpg}, \code{png}, \code{svg} and \code{pdf}. Passing several
writes the same figure in each format under a shared name, for example
\code{-{}-figure-format jpg svg} emits both \code{...\_clamped.jpg} and
\code{...\_clamped.svg}; \code{svg} and \code{pdf} are resolution-independent
vector formats suited to publications. Pass
\code{-{}-no-figures} to skip the figures (while
still writing spreadsheets), or \code{-{}-no-save} to skip all output files.

\noindent For example, to render each figure as both a JPEG (for quick preview)
and an SVG (for inclusion in a report or further editing):
\begin{lstlisting}[style=shellstyle]
python window_deformation.py --figure-format jpg svg
\end{lstlisting}
\noindent which writes, for the default single case,
\code{single\_case\_deflection\_fused\_silica\_clamped.jpg} and
\code{single\_case\_deflection\_fused\_silica\_clamped.svg} (plus the matching
stress pair).
Because the overwrite guard reserves a stem across \emph{all} requested formats
at once, re-running the command bumps every format together
(\code{..\_clamped\_1.jpg} and \code{..\_clamped\_1.svg}), so the raster and
vector copies of one figure never drift apart.

\paragraph{Spreadsheet output.} Alongside every figure the tool writes an Excel
workbook (\code{.xlsx}) containing the same numeric results echoed to the console,
into the same output directory. A single case produces
\code{single\_case\_fused\_silica\_clamped.xlsx} with two worksheets: \emph{Single Case}
(a labelled quantity/value/units summary table, whose first row names the
analysed window material) and \emph{Profile Data},
which tabulates the full radial profiles behind the plotted curves---one row
per solver sample point giving radius, transverse deflection and the four
stress components (radial and hoop stress on the top and bottom faces).
A single-parameter sweep produces
\code{sweep\_\textit{variable}\_fused\_silica\_clamped.xlsx} (one row per swept
value), prefixed with a \emph{parameters held constant} block that lists the
material, the fixed geometry and load, boundary condition, plate theory, Young's
modulus, Poisson's ratio and allowable stress (the swept variable is omitted); and the combined
sweep produces \code{sweep\_combined\_fused\_silica\_clamped.xlsx} with four worksheets. The first,
\emph{Case Details}, lists every grid case with the same per-case quantities
reported for a single case---thickness, diameter, thickness/diameter ratio,
resolved plate theory, linear and nonlinear centre deflection, peak tensile stress
and safety factor (one row per thickness/diameter combination). The remaining three
worksheets---center deflection, peak tensile stress and safety factor---each present
that response as a thickness $\times$ diameter matrix. Every combined-sweep
worksheet is prefixed with a \emph{parameters held constant} block listing the
material, pressure, boundary condition, plate theory, Young's modulus, Poisson's
ratio and allowable stress used for the grid.
Spreadsheet cells hold full-precision values displayed to five
decimals. Every workbook also carries a \emph{Note} row stating the
plate-theory selection rule---under \code{auto}, based on the radius-to-thickness
ratio $R/h = 1/(2\,t/D)$, Kirchhoff thin-plate theory is used for
$t/D < \num{0.05}$ ($R/h > 10$) and Mindlin--Reissner thick-plate theory for
$t/D \ge \num{0.05}$ ($R/h \le 10$, spanning the transition zone
$\num{0.05}\le t/D<\num{0.1}$ and the definite thick-plate region $t/D\ge\num{0.1}$),
while an explicit \code{-{}-plate-theory} choice overrides it.
Pass \code{-{}-no-excel} to skip the spreadsheets (while still writing
figures), or \code{-{}-no-save} to skip all output files. Excel export uses
\code{openpyxl}.

\paragraph{Output file naming.} Every figure and spreadsheet name embeds the
window material and always ends with the boundary condition
(\code{\textit{name}\_fused\_silica\_clamped},
\code{\textit{name}\_methyl\_methacrylate\_simply\_supported}) so results are
self-identifying and runs with different materials or edge conditions never
collide. The material token is a lowercase, filename-safe form of the selected
material's display name (for example \code{Fused Silica}
$\rightarrow$ \code{fused\_silica}, \code{N-BK7} $\rightarrow$ \code{n\_bk7},
\code{96\% Silica (Vycor)} $\rightarrow$ \code{96\_silica\_vycor}). The tool
also never overwrites an existing output file: if the material/boundary name is
already taken, a numeric suffix is added
(\code{\textit{name}\_fused\_silica\_clamped\_1},
\code{\textit{name}\_fused\_silica\_clamped\_2}, \ldots) to guarantee
uniqueness. The console always reports the exact path written.

\paragraph{Custom output names.} The default base names (\code{single\_case},
\code{sweep\_\textit{variable}}, \code{sweep\_combined}) can be overridden with
\code{-{}-figure-name} and \code{-{}-excel-name}. The material token,
boundary-condition suffix, uniqueness suffix and file extension are still
appended, so with the default fused-silica preset \code{-{}-excel-name results}
writes \code{results\_fused\_silica\_clamped.xlsx}. In single-case
mode the two figures take a plot-type suffix, so \code{-{}-figure-name mywindow}
writes \code{mywindow\_deflection\_fused\_silica\_clamped.jpg} and
\code{mywindow\_stress\_fused\_silica\_clamped.jpg}. For example,
\begin{quote}
\code{python window\_deformation.py -{}-figure-name mywindow -{}-excel-name results}
\end{quote}
names the default single case \code{mywindow\_deflection\_fused\_silica\_clamped.jpg},
\code{mywindow\_stress\_fused\_silica\_clamped.jpg} and
\code{results\_fused\_silica\_clamped.xlsx}.

\paragraph{Diameter sweep defaults.} When \code{-{}-sweep-variable diameter\_mm} is
selected without explicit \code{-{}-sweep-start}/\code{-{}-sweep-stop}, the sweep
defaults to diameters \qtyrange{250}{450}{\milli\meter} in $5$ steps (radius \qtyrange{125}{225}{\milli\meter} in
\qty{25}{\milli\meter} increments).

\paragraph{Combined sweep.} Selecting \code{-{}-sweep-variable all} evaluates a
full thickness $\times$ diameter grid at the fixed base pressure and presents it
in a single figure (\code{sweep\_combined\_fused\_silica\_clamped.jpg}) as \emph{families} of curves. The
figure is a $2\times 3$ grid of panels: the top row plots each response versus
thickness with one curve per diameter, and the bottom row plots each response
versus diameter with one curve per thickness. The three columns are center
deflection, peak tensile stress and safety factor. The thickness axis defaults to
\qtyrange{15}{35}{\milli\meter} in $5$ steps and the diameter axis to \qtyrange{250}{450}{\milli\meter} in $5$ steps
($25$ solved cases). The pressure sweep is intentionally excluded so that the
whole grid shares one loading condition. Each grid leg has its own range flags:
\code{-{}-thickness-sweep-start}/\code{-{}-thickness-sweep-stop}/\code{-{}-thickness-sweep-count}
control the thickness axis and
\code{-{}-diameter-sweep-start}/\code{-{}-diameter-sweep-stop}/\code{-{}-diameter-sweep-count}
control the diameter axis; any flag left unset falls back to the default above, so a
single leg can be widened without respecifying the other. The plain
\code{-{}-sweep-start}/\code{-{}-sweep-stop}/\code{-{}-sweep-count} values apply only to
single-variable sweeps and are ignored in this mode; if any of them is supplied
alongside \code{-{}-sweep-variable all} the program prints a boxed
\emph{ignored arguments} warning to \code{stderr} naming the offending flags and
pointing to the grid-range flags, then continues with the default (or
grid-flag) ranges. The console prints three
matrices (deflection, stress and safety factor) with thickness as rows and diameter
as columns.

\paragraph{Fixed-parameter annotations.} Every figure states which parameters are
held fixed through a compact two-line subtitle: the top line lists the analysed
material together with the fixed geometry and load, and the second line carries the
boundary condition and plate theory. The single-case deflection and stress figures
read, for example, \emph{Material = Fused Silica, Diameter = 400.0 mm,
Thickness = 25.0 mm, Pressure = 101,000 Pa} on the first line and
\emph{Edge = clamped, Theory = Mindlin (t/D = 0.062)} on the second, while the
plot titles themselves stay plain (\emph{Deflection Profile}, \emph{Stress
Profiles}). A single-parameter sweep uses the same two-line subtitle but omits the
swept variable from the fixed list. The combined figure places the globally common
constants (material, pressure, edge support and plate theory) in the two-line
suptitle; each panel's legend identifies the curve family (diameter for the
thickness panels, thickness for the diameter panels).

% ===========================================================================
\section{Examples}
% ===========================================================================

\subsection{Default single case}

\begin{lstlisting}[style=shellstyle]
python window_deformation.py
\end{lstlisting}

\noindent Produces the console summary below and writes
\code{figures/single\_case\_deflection\_fused\_silica\_clamped.jpg} and
\code{figures/single\_case\_stress\_fused\_silica\_clamped.jpg}:

\begin{lstlisting}[style=shellstyle]
=== Fused Silica Window Deformation (von Karman Nonlinear Plate BVP) ===
Diameter: 450.000 mm, Thickness: 25.000 mm, Pressure: 101000.000 Pa
Boundary condition: clamped

Linear center deflection:    0.0419 mm
Nonlinear center deflection: 0.0419 mm
Peak tensile stress:         6.137 MPa
Safety factor (allowable/peak): 0.764
Saved figure: figures\single_case_deflection_fused_silica_clamped.jpg
Saved figure: figures\single_case_stress_fused_silica_clamped.jpg
\end{lstlisting}

For this geometry $w/t \approx 0.0017$, so the membrane stiffening is negligible
and the nonlinear result coincides with linear theory --- as it physically should.

\subsection{Simply supported edge}

\begin{lstlisting}[style=shellstyle]
python window_deformation.py --boundary-condition simply_supported
\end{lstlisting}

A simply supported edge permits rotation and therefore deflects substantially more
than a clamped edge ($\approx$~\qty{0.187}{\milli\meter} centre deflection for the default
geometry).

\subsection{Shear-deformable (Mindlin) thick plate}
\label{sec:mindlin-example}

\begin{lstlisting}[style=shellstyle]
python window_deformation.py --diameter-mm 250 --thickness-mm 35 \
    --plate-theory mindlin
\end{lstlisting}

For a \qty{35}{\milli\meter} plate on a \qty{250}{\milli\meter} aperture the thickness-to-diameter ratio is
$t/D = 0.14$, beyond the thin-plate limit. The Mindlin--Reissner solution adds the
transverse-shear compliance of \eqref{eq:mindlin-clamped} and therefore reports a
larger centre deflection than the Kirchhoff model, while the peak stress is
essentially unchanged. With the default \code{-{}-plate-theory auto} the same case
selects Mindlin automatically because $t/D = 0.14 \ge 0.05$ ($R/h < 10$); passing
\code{kirchhoff} forces the thin-plate model for comparison.

\subsection{Selecting a material preset}
\label{sec:material-example}

\begin{lstlisting}[style=shellstyle]
python window_deformation.py --material sapphire
\end{lstlisting}

\noindent Solves the default geometry with the sapphire preset
(Table~\ref{tab:materials}); the console header echoes the material name
(\emph{Sapphire Window Deformation}). The preset values can be tuned on the same
command line---for example, applying a more conservative allowable stress while
keeping sapphire's stiffness:

\begin{lstlisting}[style=shellstyle]
python window_deformation.py --material sapphire --allowable-stress-mpa 200
\end{lstlisting}

\noindent To see every available preset without running a solve:

\begin{lstlisting}[style=shellstyle]
python window_deformation.py --list-materials
\end{lstlisting}

\subsection{Thickness design sweep}

\begin{lstlisting}[style=shellstyle]
python window_deformation.py --sweep-variable thickness_mm \
    --sweep-start 15 --sweep-stop 35 --sweep-count 5
\end{lstlisting}

\noindent Produces an aligned table (and writes\\
\code{figures/sweep\_thickness\_mm\_fused\_silica\_clamped.jpg}):

\begin{lstlisting}[style=shellstyle]
=== Design Sweep Summary ===
Sweep variable: Thickness (mm)

     Value    Deflection(mm)   PeakStress(MPa)    SafetyFactor
--------------------------------------------------------------
    15.000           0.19394          17.06445          0.2747
    20.000           0.08183           9.59089          0.4888
    25.000           0.04389           6.13685          0.7640
    30.000           0.02591           4.26132          1.1002
    35.000           0.01669           3.13064          1.4976
\end{lstlisting}

\subsection{Diameter design sweep}

With no explicit bounds, the diameter sweep defaults to \qtyrange{250}{450}{\milli\meter} in $5$
steps (radius \qtyrange{125}{225}{\milli\meter}):

\begin{lstlisting}[style=shellstyle]
python window_deformation.py --sweep-variable diameter_mm
\end{lstlisting}

\noindent Produces (and writes \code{figures/sweep\_diameter\_mm\_fused\_silica\_clamped.jpg}):

\begin{lstlisting}[style=shellstyle]
=== Design Sweep Summary ===
Sweep variable: Diameter (mm)

     Value    Deflection(mm)   PeakStress(MPa)    SafetyFactor
--------------------------------------------------------------
   250.000           0.00461           1.89379          2.4757
   300.000           0.00916           2.72711          1.7192
   350.000           0.01654           3.71201          1.2630
   400.000           0.02773           4.84855          0.9670
   450.000           0.04389           6.13685          0.7640
\end{lstlisting}

\subsection{Pressure design sweep}

\begin{lstlisting}[style=shellstyle]
python window_deformation.py --sweep-variable pressure_pa \
    --sweep-start 50000 --sweep-stop 200000 --sweep-count 6
\end{lstlisting}

\subsection{Combined thickness and diameter sweep}

To study both geometry parameters at once under the fixed base pressure, select
\code{all}. Three response matrices are printed (deflection, peak stress and
safety factor) and a single combined figure of curve families
(\code{figures/sweep\_combined\_fused\_silica\_clamped.jpg}) is written:

\begin{lstlisting}[style=shellstyle]
python window_deformation.py --sweep-variable all
\end{lstlisting}

\noindent Console output (abbreviated):

\begin{lstlisting}[style=shellstyle, basicstyle=\ttfamily\footnotesize]
=== Combined Sweep: Center Deflection (mm) ===
Rows: Thickness (mm)  |  Columns: Diameter (mm)
   Thk\Dia         250.0         300.0         350.0         400.0         450.0
--------------------------------------------------------------------------------
      15.0       0.01950       0.03979       0.07098       0.12108       0.19394
      ...
      35.0       0.00189       0.00365       0.00645       0.01066       0.01669

=== Combined Sweep: Peak Tensile Stress (MPa) ===
   ... (thickness x diameter matrix) ...

=== Combined Sweep: Safety Factor ===
   ... (thickness x diameter matrix) ...
\end{lstlisting}

\noindent Both grid legs accept their own range flags. For example, to sweep the
thickness from \qtyrange{25}{70}{\milli\meter} in $10$ steps on a simply supported
edge while leaving the diameter leg at its default (\qtyrange{250}{450}{\milli\meter}):

\begin{lstlisting}[style=shellstyle]
python window_deformation.py --boundary-condition simply_supported \
    --sweep-variable all \
    --thickness-sweep-start 25 --thickness-sweep-stop 70 \
    --thickness-sweep-count 10
\end{lstlisting}

\noindent Either leg can be overridden independently. This widens \emph{both} axes at
once (thickness \qtyrange{25}{70}{\milli\meter} in $10$ steps and diameter
\qtyrange{200}{500}{\milli\meter} in $7$ steps):

\begin{lstlisting}[style=shellstyle]
python window_deformation.py --sweep-variable all \
    --thickness-sweep-start 25 --thickness-sweep-stop 70 --thickness-sweep-count 10 \
    --diameter-sweep-start 200 --diameter-sweep-stop 500 --diameter-sweep-count 7
\end{lstlisting}

\noindent Any flag left unset keeps its default, so a single leg can be retuned
without respecifying the other. The plain \code{-{}-sweep-start}/\code{-{}-sweep-stop}/\code{-{}-sweep-count}
flags are for single-variable sweeps and have no effect in \code{all} mode; supplying
one of them here triggers a boxed warning that names the ignored flag(s), for example:

\begin{lstlisting}[style=shellstyle, basicstyle=\ttfamily\footnotesize]
$ python window_deformation.py --sweep-variable all --sweep-start 10
****************************************************************************
*  !!! IGNORED ARGUMENTS WARNING !!!                                       *
*                                                                          *
*  Ignoring --sweep-start in combined ('all') sweep mode. These flags      *
*  apply only to a single-variable sweep. To set the grid ranges use       *
*  --thickness-sweep-start/-stop/-count and                                *
*  --diameter-sweep-start/-stop/-count instead.                            *
****************************************************************************
\end{lstlisting}

\subsection{Programmatic use}

The module can be imported and used directly, which is convenient for scripting or
notebooks:

\begin{lstlisting}[style=pystyle]
import window_deformation as wd

config = wd.PlateConfig.from_engineering_units(
    diameter_mm=450.0,
    thickness_mm=25.0,
    pressure_pa=101_000.0,
    boundary_condition=wd.CLAMPED,
)
solution = wd.solve_plate(config, n_points=500)

print(solution.center_deflection_mm)  # mm
print(solution.max_tensile_mpa)       # MPa
print(solution.safety_factor)
\end{lstlisting}

\subsection{Custom material}

\begin{lstlisting}[style=pystyle]
import window_deformation as wd

bk7 = wd.Material(
    youngs_modulus_pa=82e9,
    poisson_ratio=0.21,
    allowable_stress_pa=40e6,
    name="N-BK7",
)
config = wd.PlateConfig.from_engineering_units(
    diameter_mm=300.0, thickness_mm=20.0, pressure_pa=101_000.0, material=bk7,
)
solution = wd.solve_plate(config)
\end{lstlisting}

\noindent The named presets exposed on the command line are also available
programmatically through the \code{wd.MATERIAL\_PRESETS} dictionary (keyed by the
same names as \code{-{}-material}), so \code{wd.MATERIAL\_PRESETS["sapphire"]}
returns a ready-made \code{Material}.

% ===========================================================================
\section{Output Quantities}
% ===========================================================================

\begin{description}[style=nextline]
  \item[Linear center deflection] $w_{0}^{\text{lin}}$ from
    Eq.~\eqref{eq:clamped-linear} or~\eqref{eq:ss-linear}: the small-deflection
    (bending-only) baseline. It scales in direct proportion to pressure and
    ignores membrane stretching, so it is a quick sanity check rather than the
    design value.
  \item[Nonlinear center deflection] $w(0)$ from the full von K\'arm\'an solve,
    which adds the membrane stiffening that develops as the mid-surface stretches.
    This is the physically complete value and the one recommended for design.
  \item[Peak tensile stress] $\sigma_{\max}$, the governing quantity for brittle
    fracture.
  \item[Safety factor] $\sigma_{\text{allow}} / \sigma_{\max}$.
\end{description}

\paragraph{Which deflection to use.} Report and design with the \emph{nonlinear}
center deflection; it is the value carried through every other output (the plotted
curves, the combined-grid matrices and the stress/safety results are all consistent
with it), while the linear column is provided only as a reference. The two agree
closely while the normalised deflection $w_{0}/t$ (center deflection divided by
thickness) is small: for $w_{0}/t \lesssim 0.2$--$0.3$ the plate is effectively in
the linear-bending regime and either number is fine. As $w_{0}/t$ grows, membrane
stiffening makes the nonlinear value increasingly smaller than the linear one, so
the gap between the two columns is a direct measure of how much the large-deflection
correction matters. Above $w_{0}/t = \num{5}$ (\code{VON\_KARMAN\_WT\_WARN\_RATIO})
even the von K\'arm\'an model is out of its moderate-rotation validity envelope; the
tool then prints a boxed model-validity warning (and highlights the affected
spreadsheet rows) advising a large-rotation shell or finite-element model. Every
spreadsheet also carries an italic note next to the results restating this guidance.
Note that when the shear-deformable Mindlin theory is selected (thick plates, low
$w_{0}/t$) the nonlinear value can be slightly \emph{larger} than the pure-bending
linear value, because transverse-shear softening then outweighs the small membrane
stiffening.

The tool writes its figures as JPEG files by default (no interactive window), and
can additionally or alternatively emit PNG, SVG or PDF via \code{-{}-figure-format}.
A single case produces a deflection profile and a stress profile (radial and hoop
on both faces); a sweep produces a two-panel figure of deflection versus the swept
variable and peak-stress/safety-factor versus the swept variable. Files are named
descriptively and written to the output directory (\code{figures/} by default).

% ===========================================================================
\section{Validation and Testing}
% ===========================================================================

The package ships with a \code{pytest} suite (\code{test\_window\_deformation.py})
of 37 tests covering:

\begin{itemize}[noitemsep]
  \item input validation and edge cases (non-positive dimensions, out-of-range
    Poisson ratio, invalid boundary condition, degenerate mesh sizes);
  \item physical correctness (zero edge deflection, centre is the maximum,
    monotonic profile, clamped versus simply-supported ordering, thickness and
    pressure trends);
  \item agreement of the nonlinear solver with linear theory in the thin-deflection
    limit and against the closed-form expressions;
  \item safety-factor scaling and design-sweep behaviour; and
  \item command-line plumbing.
\end{itemize}

Run the tests and the linter with:

\begin{lstlisting}[style=shellstyle]
python -m pytest
python -m ruff check .
\end{lstlisting}

% ===========================================================================
\section{Modelling Assumptions and Limitations}
% ===========================================================================

\begin{itemize}
  \item \textbf{Plate kinematics.} By default \code{auto} selects the plate model
    from the radius-to-thickness ratio: Kirchhoff--Love kinematics (no transverse
    shear) for genuinely thin plates ($t/D < 0.05$, $R/h > 10$), and the
    shear-deformable Mindlin--Reissner model (Section~\ref{sec:mindlin}) for the
    transition zone and thicker ($t/D \ge 0.05$, $R/h \le 10$). The default window
    has $t/D = 0.056$ ($R/h = 9$), which lies in the transition zone, so
    \code{auto} resolves it to Mindlin; the \code{-{}-plate-theory} flag can force
    either model explicitly.
  \item \textbf{Linear-elastic, isotropic material.} No plasticity, viscoelasticity,
    or anisotropy is modelled.
  \item \textbf{Radially constrained edge.} The edge is fixed against in-plane
    radial motion ($u(a)=0$). A free-to-slide edge would reduce membrane
    stiffening.
  \item \textbf{Uniform static pressure.} The load is a uniform, static pressure on
    one face. Dynamic, thermal, and self-weight effects are not included.
  \item \textbf{Brittle failure criterion.} The safety factor uses a maximum
    tensile stress criterion appropriate for brittle materials such as fused
    silica; it does not account for statistical strength (e.g.\ Weibull) or
    surface-flaw distributions.
\end{itemize}

% ===========================================================================
\section{References}
% ===========================================================================

\begin{enumerate}[noitemsep]
  \item S.~Timoshenko and S.~Woinowsky-Krieger, \emph{Theory of Plates and
    Shells}, 2nd ed., McGraw-Hill, 1959.
  \item L.~D.~Landau and E.~M.~Lifshitz, \emph{Theory of Elasticity},
    3rd ed., Pergamon Press, 1986 (F\"oppl--von K\'arm\'an equations).
  \item P.~Virtanen \emph{et al.}, ``SciPy 1.0: Fundamental Algorithms for
    Scientific Computing in Python,'' \emph{Nature Methods}, 17, 261--272, 2020
    (the \code{solve\_bvp} collocation solver).
  \item Brookhaven National Laboratory, ``Glass and Plastic Window Design for
    Pressure Vessels,'' \S1.4.2 of the \emph{Occupational Health and Safety
    Guide} (Interim), Brookhaven National Laboratory, Upton, NY. Source of the
    material properties, allowable design stresses and table Note~\#8 used for
    the built-in presets.\label{ref:bnl}
\end{enumerate}

\end{document}
